\begin{derivation}[从\eqnrefs{eq:lorentz-irrep-dictionary-dp8}到\eqnrefs{eq:weyl-vector-product-r68}]
只写维数 $2\times2=4$ 不能证明一个双旋量就是四矢量；关键是它的变换律必须与正文已经建立的
$X(x)\mapsto AX(x)A^\dagger$ 完全相同。设左手 Weyl 旋量按
\begin{equation*}
 \xi_\alpha\longmapsto \xi'_\alpha
 =A_\alpha{}^\beta\xi_\beta,
 \qquad A\in SL(2,\mathbb C),
\end{equation*}
变换。其厄米共轭带点指标按
\begin{equation*}
 \bar\chi_{\dot\alpha}\longmapsto
 \bar\chi'_{\dot\alpha}
 =(A^\dagger)_{\dot\alpha}{}^{\dot\beta}
 \bar\chi_{\dot\beta}
\end{equation*}
变换。构造一个左右各带一个指标的对象
\begin{equation*}
 V_{\alpha\dot\beta}=\xi_\alpha\bar\chi_{\dot\beta}.
\end{equation*}
把 $V_{\alpha\dot\beta}$ 看成 $2\times2$ 矩阵，则
\begin{align*}
 V'_{\alpha\dot\beta}
 &=A_\alpha{}^\gamma
 V_{\gamma\dot\delta}
 (A^\dagger)^{\dot\delta}{}_{\dot\beta},\\
 \Longleftrightarrow\qquad
 V'&=AVA^\dagger.
\end{align*}
这已经与正文\eqnrefs{eq:sl2c-action-X-dp8} 的四矢量矩阵变换完全相同。

为了把矩阵分量显式还原成四矢量，使用
$\sigma^\mu=(I_2,\boldsymbol\sigma_{\rm P})$、
$\bar\sigma^\mu=(I_2,-\boldsymbol\sigma_{\rm P})$。Weyl 交织恒等式给
\begin{equation*}
 \operatorname{Tr}(\bar\sigma^\mu\sigma^\nu)=2\eta^{\mu\nu}.
\end{equation*}
因此展开
\begin{equation*}
 V=V_\nu\sigma^\nu
\end{equation*}
的逆变换不是一个符号约定，而由迹正交性唯一固定：
\begin{equation*}
 V^\mu=\frac12\operatorname{Tr}(\bar\sigma^\mu V).
\end{equation*}
对变换后的矩阵，
\begin{align*}
 V'^\mu
 &=\frac12\operatorname{Tr}(\bar\sigma^\mu A\sigma^\nu A^\dagger)V_\nu\\
 &\equiv\Lambda^\mu{}_{\nu}(A)V^\nu .
\end{align*}
这里
\begin{equation*}
 \Lambda^\mu{}_{\nu}(A)
 =\frac12\operatorname{Tr}
 \left(\bar\sigma^\mu A\sigma_\nu A^\dagger\right)
\end{equation*}
正是 $SL(2,\mathbb C)$ 双覆盖诱导出的 Lorentz 矩阵。进一步，
\begin{equation*}
 \det V=V_\mu V^\mu,
 \qquad
 \det(AVA^\dagger)=|\det A|^2\det V=\det V,
\end{equation*}
所以
$V'_\mu V'^\mu=V_\mu V^\mu$；这再次验证 $\Lambda(A)$ 保持 Minkowski 度规。

因此
\begin{equation*}
 V_{\alpha\dot\beta}=V_\mu(\sigma^\mu)_{\alpha\dot\beta}
\end{equation*}
不是靠“维数碰巧都是四”猜出的字典，而是因为一左一右 Weyl 指标的矩阵变换律与四矢量的 $SL(2,\mathbb C)$ 实现逐项相同。这才证明
$(\tfrac12,0)\otimes(0,\tfrac12)=(\tfrac12,\tfrac12)$，即正文\eqnrefs{eq:weyl-vector-product-r68}。
\end{derivation}

\begin{derivation}[从\eqnrefs{eq:weyl-same-chirality-product-r68}到\eqnrefs{eq:twoform-chiral-decomposition-r68}]
对两个左手 Weyl 指标，任意二阶对象 $T_{\alpha\beta}$ 都有唯一的对称--反对称分解
\begin{equation*}
 T_{\alpha\beta}=T_{(\alpha\beta)}+T_{[\alpha\beta]}.
\end{equation*}
二维旋量空间中反对称二阶张量只有一个独立分量，因此必可写成
\begin{equation*}
 T_{[\alpha\beta]}=\phi\,\epsilon_{\alpha\beta}.
\end{equation*}
为什么这个系数 $\phi$ 是 Lorentz 标量？因为 $A\in SL(2,\mathbb C)$ 满足
\begin{equation*}
 A_\alpha{}^\gamma A_\beta{}^\delta
 \epsilon_{\gamma\delta}
 =(\det A)\epsilon_{\alpha\beta}
 =\epsilon_{\alpha\beta}.
\end{equation*}
所以反对称的一维子空间保持不变，正是 $(0,0)$。剩下的对称矩阵满足
$T_{(11)},T_{(12)},T_{(22)}$ 三个独立复分量，并在
\begin{equation*}
 T_{(\alpha\beta)}\longmapsto
 A_\alpha{}^\gamma A_\beta{}^\delta
 T_{(\gamma\delta)}
\end{equation*}
下封闭，因此构成三维不可约表示 $(1,0)$。于是从实际指标分解得到
\begin{equation*}
 (\tfrac12,0)\otimes(\tfrac12,0)
 =(0,0)\oplus(1,0).
\end{equation*}
右手带点指标完全平行，给出
\begin{equation*}
 (0,\tfrac12)\otimes(0,\tfrac12)
 =(0,0)\oplus(0,1).
\end{equation*}

现在转向四维反对称张量 $F_{\mu\nu}=-F_{\nu\mu}$。它有
$\binom42=6$ 个独立分量。复化以后定义 Hodge 算符
\begin{equation*}
 (\star F)^{\mu\nu}
 =\frac12\epsilon^{\mu\nu\rho\sigma}F_{\rho\sigma}.
\end{equation*}
在当前 Lorentz 度规和 $\epsilon^{0123}=+1$ 约定下，二形式满足
$\star^2=-I$，所以其本征值是 $\pm\ii$。相应投影算符为
\begin{equation*}
 \mathcal P_\pm=\frac12(I\mp\ii\star),
 \qquad
 F_\pm\equiv\mathcal P_\pm F,
 \qquad
 \star F_\pm=\pm\ii F_\pm.
\end{equation*}
因为
$\mathcal P_+\mathcal P_-=0$、
$\mathcal P_++\mathcal P_-=I$，六维复二形式空间严格分成两个互补的三维子空间：
\begin{equation*}
 \Lambda^2_\mathbb C
 =\Lambda^2_+\oplus\Lambda^2_-.
\end{equation*}
Lorentz 变换与 Hodge 算符对易，所以两个三维子空间各自保持不变。与上面同手征双旋量的两个三维不可约表示比较，分别对应 $(1,0)$ 与 $(0,1)$；因此
\begin{equation*}
 F^{\mu\nu}\in(1,0)\oplus(0,1),
\end{equation*}
即正文\eqnrefs{eq:twoform-chiral-decomposition-r68}。对电磁场，$F_\pm$ 的三个独立复分量可等价组织成 $\mathbf E\pm\ii c\mathbf B$；这不是另外引入的场，而是同一个六分量 $F_{\mu\nu}$ 在 Hodge 本征子空间中的两种手征投影。
\end{derivation}

\begin{derivation}[从\eqnrefs{eq:clifford-general}到\eqnrefs{eq:trace-four-gamma-r9}]

未偏振费米子振幅在完成自旋求和后会变成 Dirac 矩阵迹。正文已经把 Clifford 代数固定为\eqnrefs{eq:clifford-general}；这里仅用反对易关系与迹循环性，把它化成正文\eqnrefs{eq:trace-two-gamma-r9,eq:trace-four-gamma-r9}，从而得到后续树级截面最常用的两个迹公式。



由已经建立的关系
\begin{equation*}
 \{\gamma^\mu,\gamma^\nu\}=2\eta^{\mu\nu}\id,
 \qquad
 \Tr\id=4.
\end{equation*}
对 Clifford 关系取迹：
\begin{align*}
 \Tr(\gamma^\mu\gamma^\nu)
 +\Tr(\gamma^\nu\gamma^\mu)
 &=2\eta^{\mu\nu}\Tr\id,\\
 2\Tr(\gamma^\mu\gamma^\nu)
 &=8\eta^{\mu\nu},
\end{align*}
其中第二行使用迹的循环性。故
\begin{equation*}
 {\Tr(\gamma^\mu\gamma^\nu)=4\eta^{\mu\nu}.}
 \end{equation*}
对四个 gamma，先把最左边 $\gamma^\mu$ 逐次移过后面的矩阵：
\begin{align*}
 T^{\mu\nu\rho\sigma}
 &\equiv\Tr(\gamma^\mu\gamma^\nu\gamma^\rho\gamma^\sigma),\\
 &=\Tr[(2\eta^{\mu\nu}-\gamma^\nu\gamma^\mu)
 \gamma^\rho\gamma^\sigma],\\
 &=2\eta^{\mu\nu}\Tr(\gamma^\rho\gamma^\sigma)
 -\Tr(\gamma^\nu\gamma^\mu\gamma^\rho\gamma^\sigma).
\end{align*}
再在第二项中把 $\gamma^\mu$ 移过 $\gamma^\rho$，
\begin{align*}
 T^{\mu\nu\rho\sigma}
 &=8\eta^{\mu\nu}\eta^{\rho\sigma}
 -2\eta^{\mu\rho}\Tr(\gamma^\nu\gamma^\sigma)
 +\Tr(\gamma^\nu\gamma^\rho\gamma^\mu\gamma^\sigma).
\end{align*}
最后用迹的循环性
\begin{equation*}
 \Tr(\gamma^\nu\gamma^\rho\gamma^\mu\gamma^\sigma)
 =\Tr(\gamma^\mu\gamma^\sigma\gamma^\nu\gamma^\rho)
\end{equation*}
并重复一次反对易，得到
\begin{equation*}
 {
 \Tr(\gamma^\mu\gamma^\nu\gamma^\rho\gamma^\sigma)
 =4(\eta^{\mu\nu}\eta^{\rho\sigma}
 -\eta^{\mu\rho}\eta^{\nu\sigma}
 +\eta^{\mu\sigma}\eta^{\nu\rho}).}
 \end{equation*}
这条式子是后面所有未偏振树级费米子截面的核心；后续未偏振费米子散射的自旋求和将直接使用这一由 Clifford 代数固定的结果。
\end{derivation}

\begin{derivation}[从\eqnrefs{eq:clifford-general}到\eqnrefs{eq:trace-even-recursion-r68}]
正文\eqnrefs{eq:trace-even-recursion-r68} 已经把 $T_{2n}^{\mu_1\cdots\mu_{2n}}$ 记为长度 $2n$ 的 Dirac 迹。把最左端 $\gamma^{\mu_1}$ 逐次移到右端。第一次交换给

\begin{equation*}
 \gamma^{\mu_1}\gamma^{\mu_2}
 =2\eta^{\mu_1\mu_2}-\gamma^{\mu_2}\gamma^{\mu_1},
\end{equation*}
第二次再对残余的 $\gamma^{\mu_1}\gamma^{\mu_3}$ 使用同一关系。一直进行到第 $2n$ 个矩阵，得到
\begin{align*}
 \gamma^{\mu_1}\cdots\gamma^{\mu_{2n}}
 =&\sum_{j=2}^{2n}(-1)^j 2\eta^{\mu_1\mu_j}
 \gamma^{\mu_2}\cdots\widehat{\gamma^{\mu_j}}\cdots\gamma^{\mu_{2n}}
 \nonumber\\
 &-\gamma^{\mu_2}\cdots\gamma^{\mu_{2n}}\gamma^{\mu_1},
 \end{align*}
其中帽号表示删除该因子；最后一项前的负号来自一共经过 $2n-1$ 次交换。取迹并用迹的循环性，最后一项等于 $-T_{2n}$，因此移到左边得到 $2T_{2n}$。两边除以 $2$：
\begin{equation*}
 {
 T_{2n}^{\mu_1\cdots\mu_{2n}}
 =\sum_{j=2}^{2n}(-1)^j
 \eta^{\mu_1\mu_j}
 \Tr\!\left(
 \gamma^{\mu_2}\cdots\widehat{\gamma^{\mu_j}}\cdots\gamma^{\mu_{2n}}
 \right).}
 \end{equation*}
$n=2$ 时立刻恢复四伽马矩阵迹；$n=3$ 时得到六个 $\gamma$ 矩阵的迹的十五个度规配对。Compton、M\o ller、Bhabha 与圈积分分子中的长矩阵链都可按这一递推逐级缩短。
\end{derivation}



\begin{derivation}[从\eqnrefs{eq:clifford-general}到\eqnrefs{eq:gamma-d-sandwich-r68}]

第一章的维数正规化已经采用对复维数 $d$ 作解析延拓的约定。若后续积分采用这一调节方法，矩阵代数也必须先保持在同一个 $d$ 维参数下；这里从正文 Clifford 母式直接缩并一个指标，得到正文\eqnrefs{eq:gamma-d-sandwich-r68}。



由 Clifford 代数
\begin{align*}
 \gamma^\mu\gamma^\alpha\gamma_\mu
 &=(2\eta^{\mu\alpha}-\gamma^\alpha\gamma^\mu)\gamma_\mu,\\
 &=2\gamma^\alpha-\gamma^\alpha(\gamma^\mu\gamma_\mu).
\end{align*}
在 $d$ 维中
\begin{equation*}
 \gamma^\mu\gamma_\mu=d\id,
\end{equation*}
因此
\begin{equation*}
 {
 \gamma^\mu\gamma^\alpha\gamma_\mu=(2-d)\gamma^\alpha.}
 \end{equation*}
四维时就是 $-2\gamma^\alpha$。自能分子的约化直接使用这一恒等式。
\end{derivation}

\begin{derivation}[从\eqnrefs{eq:clifford-general}到\eqnrefs{eq:gamma-three-sandwich-r10}]
M\o ller、Bhabha、Compton 与圈分子中经常出现把一串 $\gamma$ 矩阵夹在 $\gamma^\mu$ 与 $\gamma_\mu$ 之间的结构。它们都可逐次用

\begin{equation*}
 \gamma^\mu\gamma^\nu=2\eta^{\mu\nu}-\gamma^\nu\gamma^\mu
 \end{equation*}
把左端的 $\gamma^\mu$ 向右移动得到。四维中首先有
\begin{align*}
 \gamma^\mu\slashed a\gamma_\mu
 &=a_\alpha\gamma^\mu\gamma^\alpha\gamma_\mu,\\
 &=a_\alpha(2\gamma^\alpha-\gamma^\alpha\gamma^\mu\gamma_\mu),\\
 &=(2-4)\slashed a=-2\slashed a.
 \end{align*}
对两个斜线记号，
\begin{align*}
 \gamma^\mu\slashed a\slashed b\gamma_\mu
 &=\gamma^\mu\slashed a(2b_\mu-\gamma_\mu\slashed b),\\
 &=2\slashed b\slashed a-
 (\gamma^\mu\slashed a\gamma_\mu)\slashed b,\\
 &=2\slashed b\slashed a+2\slashed a\slashed b,\\
 &=2\{\slashed a,\slashed b\}=4a\!\cdot b\,\id_4.
 \end{align*}
对三个斜线记号，先把最右的 $\slashed c$ 留在外面：
\begin{align*}
 \gamma^\mu\slashed a\slashed b\slashed c\gamma_\mu
 &=2\slashed c\slashed a\slashed b
 -\gamma^\mu\slashed a\slashed b\gamma_\mu\slashed c,\\
 &=2\slashed c\slashed a\slashed b-4(a\!\cdot b)\slashed c.
\end{align*}
再把第一项中的斜线记号重新排序，得到等价形式
\begin{equation*}
 {
 \gamma^\mu\slashed a\slashed b\slashed c\gamma_\mu
 =-2\slashed c\slashed b\slashed a.}
 \end{equation*}
把 $a,b,c$ 换成实际外腿四动量，就得到后续干涉迹中使用的短矩阵链化简式。
\end{derivation}

\begin{derivation}[从\eqnrefs{eq:gamma5-def}到\eqnrefs{eq:gamma5-four-trace-r68}]

轴矢量流和手征结构需要一个与普通四伽马迹不同的对象。正文\eqnrefs{eq:gamma5-def} 已固定 $\gamma^5$ 与取向约定；先证明相应秩四迹完全反对称，再用 $0123$ 分量固定系数，从而得到正文\eqnrefs{eq:gamma5-four-trace-r68}。



固定
\begin{equation*}
 \gamma^5=\ii\gamma^0\gamma^1\gamma^2\gamma^3,
 \qquad \epsilon^{0123}=+1.
\end{equation*}
张量
\begin{equation*}
 T^{\mu\nu\rho\sigma}
 \equiv\Tr(\gamma^5\gamma^\mu\gamma^\nu\gamma^\rho\gamma^\sigma)
\end{equation*}
先证明 $T^{\mu\nu\rho\sigma}$ 的完全反对称性。交换相邻的 $\mu,\nu$ 两个指标时，Clifford 代数给出
\begin{align*}
T^{\nu\mu\rho\sigma}
&=\Tr\!\left(\gamma^5\gamma^\nu\gamma^\mu\gamma^\rho\gamma^\sigma\right),\\
&=-\Tr\!\left(\gamma^5\gamma^\mu\gamma^\nu\gamma^\rho\gamma^\sigma\right)
+2\eta^{\mu\nu}\Tr\!\left(\gamma^5\gamma^\rho\gamma^\sigma\right).
\end{align*}
第二项为零。为看清这一点，利用 $\{\gamma^5,\gamma^\rho\}=0$ 和迹的循环性：
\begin{align*}
\Tr(\gamma^5\gamma^\rho\gamma^\sigma)
&=-\Tr(\gamma^\rho\gamma^5\gamma^\sigma),\\
&=-\Tr(\gamma^5\gamma^\sigma\gamma^\rho),\\
&=-\Tr\!\left[\gamma^5(2\eta^{\rho\sigma}-\gamma^\rho\gamma^\sigma)\right],\\
&=-2\eta^{\rho\sigma}\Tr(\gamma^5)+\Tr(\gamma^5\gamma^\rho\gamma^\sigma).
\end{align*}
而 $\Tr\gamma^5=0$，所以这个两 $\gamma$ 迹为零。于是
$T^{\nu\mu\rho\sigma}=-T^{\mu\nu\rho\sigma}$。对其余任意相邻指标交换重复同一计算，得到 $T$ 对四个 Lorentz 指标完全反对称。四维中任何秩四完全反对称张量都与 $\epsilon^{\mu\nu\rho\sigma}$ 成正比，因此写成
\begin{equation*}
T^{\mu\nu\rho\sigma}=C_5\epsilon^{\mu\nu\rho\sigma}.
\end{equation*}
比例常数 $C_5$ 由 $0123$ 分量确定：
\begin{align*}
 T^{0123}
 &=\Tr(\gamma^5\gamma^0\gamma^1\gamma^2\gamma^3),\\
 &=\Tr\left[
 \ii(\gamma^0\gamma^1\gamma^2\gamma^3)^2
 \right].
\end{align*}
把第二组 $\gamma$ 矩阵逐次穿过第一组并使用
$(\gamma^0)^2=+1$、$(\gamma^i)^2=-1$，得到
\begin{equation*}
 (\gamma^0\gamma^1\gamma^2\gamma^3)^2=-\id_4,
\end{equation*}
所以
\begin{equation*}
 T^{0123}=-4\ii.
\end{equation*}
因此在上述度规与 $\epsilon^{0123}$ 约定下
\begin{equation*}
 {
 \Tr(\gamma^5\gamma^\mu\gamma^\nu\gamma^\rho\gamma^\sigma)
 =-4\ii\epsilon^{\mu\nu\rho\sigma}.}
 \end{equation*}
若改变度规或 $\epsilon^{0123}$ 的符号，右端也必须同步改变；这就是为何约定必须在全书开头冻结。
\end{derivation}

\begin{derivation}[从\eqnrefs{eq:dirac-bilinear-basis}到\eqnrefs{eq:dirac-matrix-completeness-r10}]
任意复 $4\times4$ 矩阵形成 $16$ 维线性空间。由 Clifford 代数构造

\begin{equation*}
 \mathcal B=
 \{\id,\gamma^\mu,\sigma^{\mu\nu},\gamma^\mu\gamma^5,\gamma^5\},
 \end{equation*}
其独立分量数为
$1+4+6+4+1=16$。迹恒等式说明不同 Lorentz 类型之间在迹内积
\begin{equation*}
 (A,B)\equiv\frac14\Tr(AB)
\end{equation*}
下彼此正交；例如 $\Tr\gamma^\mu=0$、$\Tr(\gamma^\mu\gamma^5)=0$、$\Tr\sigma^{\mu\nu}=0$。因此任意矩阵 $M$ 都可唯一展开为
\begin{equation*}
 M=\frac14\sum_A \Gamma_A\Tr(\Gamma^A M),
 \end{equation*}
其中 $A$ 遍历上述 16 个基元素，$\Gamma^A$ 表示按迹度规取对偶基。把 $M$ 取成只有一个元素非零的矩阵单位 $(E_{lk})_{ij}=\delta_{il}\delta_{jk}$，有
\begin{align*}
 \delta_{il}\delta_{jk}
 &=\frac14\sum_A(\Gamma_A)_{ij}
 \Tr(\Gamma^A E_{lk}),\\
 &=\frac14\sum_A(\Gamma_A)_{ij}(\Gamma^A)_{kl}.
\end{align*}
所以得到矩阵指标层面的完备恒等式
\begin{equation*}
 {
 \frac14\sum_A
 (\Gamma_A)_{ij}(\Gamma^A)_{kl}
 =\delta_{il}\delta_{kj}.}
 \end{equation*}
Fierz 重排来自完备矩阵基上的分量展开；它把
\eqnrefs{eq:dirac-matrix-completeness-r10}
插到两个旋量双线性型之间并重新配对旋量指标。若需要某一个具体 Fierz 关系，将从这条母完备关系逐项投影，而不直接引用公式表。
\end{derivation}



\begin{derivation}[从\eqnrefs{eq:sigma-barsigma-foundation-r10}到\eqnrefs{eq:weyl-determinant-mass-shell-dp21}]

由正文\eqnrefs{eq:sigma-barsigma-foundation-r10}，四动量被写成 $p\cdot\sigma$ 与 $p\cdot\bar\sigma$ 两个厄米矩阵。这里求它们的两个本征值并相乘，证明行列式正是 Lorentz 不变量 $p^2$，从而得到正文\eqnrefs{eq:weyl-determinant-mass-shell-dp21} 的质量壳条件。



写
\begin{equation*}
 p_\mu\sigma^\mu=p^0I-\mathbf p\cdot\boldsymbol\sigma.
\end{equation*}
对任意 $2\times2$ 矩阵 $M=aI+\mathbf b\cdot\boldsymbol\sigma$，Pauli 代数给
\begin{equation*}
 (\mathbf b\cdot\boldsymbol\sigma)^2=\mathbf b^2I.
\end{equation*}
因此 $\mathbf b\cdot\boldsymbol\sigma$ 的两个本征值是 $\pm|\mathbf b|$，而 $M$ 的本征值是 $a\pm|\mathbf b|$。行列式是本征值乘积：
\begin{equation*}
 \det M=(a+|\mathbf b|)(a-|\mathbf b|)=a^2-\mathbf b^2.
\end{equation*}
取 $a=p^0$、$\mathbf b=-\mathbf p$，得到
\begin{equation*}
 \det(p_\mu\sigma^\mu)
 =(p^0)^2-\mathbf p^2=p_\mu p^\mu.
\end{equation*}
对 $p_\mu\bar\sigma^\mu=p^0I+\mathbf p\cdot\boldsymbol\sigma$ 只把 $\mathbf b$ 改号，行列式不变。

质量壳 $p^2=m^2c^2$ 因而等价于
\begin{equation*}
 \det(p\cdot\sigma)=m^2c^2.
\end{equation*}
这解释了为何有质量 Weyl 方程中 $p\cdot\sigma$ 与 $p\cdot\bar\sigma$ 的乘积恢复二阶质量壳算符 $p^2-m^2c^2$；对电子而言，一阶 Dirac 方程仍是基础动力学方程。
\end{derivation}


\begin{derivation}[从\eqnrefs{eq:dirac-matrix-completeness-r10}到\eqnrefs{eq:fierz-scalar-r68}]

正文\eqnrefs{eq:dirac-matrix-completeness-r10} 是 $4\times4$ 矩阵空间的恒等分解。把它插入四个外线旋量的指标之间并改变配对顺序，就能把 $(1,2)(3,4)$ 的标量双线性改写为 $(1,4)(3,2)$ 的五类 Dirac 双线性，得到正文\eqnrefs{eq:fierz-scalar-r68}。



从矩阵完备关系
\begin{equation*}
\delta_{il}\delta_{kj}
=\frac14\sum_A(\Gamma_A)_{ij}(\Gamma^A)_{kl}
\end{equation*}
出发。对四个普通复数值 Dirac 旋量 $u_1,u_2,u_3,u_4$，标量双线性乘积写成显式矩阵指标
\begin{equation*}
(\bar u_1u_2)(\bar u_3u_4)
=(\bar u_1)_i(u_2)_j(\bar u_3)_k(u_4)_l\,
\delta_{ij}\delta_{kl}.
\end{equation*}
配对从 $(1,2)(3,4)$ 交换为 $(1,4)(3,2)$ 时， $\delta_{ij}\delta_{kl}$ 写成 $i$ 与 $l$、$k$ 与 $j$ 相连的基展开。由完备式交换哑指标得
\begin{equation*}
\delta_{ij}\delta_{kl}
=\frac14\sum_A(\Gamma_A)_{il}(\Gamma^A)_{kj}.
\end{equation*}
代入后
\begin{align*}
(\bar u_1u_2)(\bar u_3u_4)
&=\frac14\sum_A
[(\bar u_1)_i(\Gamma_A)_{il}(u_4)_l]
[(\bar u_3)_k(\Gamma^A)_{kj}(u_2)_j]\\
&=\frac14\sum_A
(\bar u_1\Gamma_Au_4)(\bar u_3\Gamma^Au_2).
\end{align*}
现在把 16 维基分成标量、矢量、张量、轴矢量、赝标量五类。按
\begin{equation*}
\sigma^{\mu\nu}=\frac{i}{2}[\gamma^\mu,\gamma^\nu]
\end{equation*}
的归一化，对偶基在张量扇区带因子 $1/2$；轴矢量扇区的迹度规带负号。于是
\begin{align*}
(\bar u_1u_2)(\bar u_3u_4)
=\frac14\Big[&
(\bar u_1u_4)(\bar u_3u_2)
+(\bar u_1\gamma^\mu u_4)(\bar u_3\gamma_\mu u_2)\\
&+\frac12(\bar u_1\sigma^{\mu\nu}u_4)
(\bar u_3\sigma_{\mu\nu}u_2)\\
&-(\bar u_1\gamma^\mu\gamma^5u_4)
(\bar u_3\gamma_\mu\gamma^5u_2)\\
&+(\bar u_1\gamma^5u_4)(\bar u_3\gamma^5u_2)
\Big].
\end{align*}
若四个对象是 Grassmann 量子场，重排 $\psi_2$ 与 $\bar\psi_3$ 到新配对顺序还会产生反对易负号。Fierz 系数由矩阵完备性决定，统计符号则由场算符次序决定。
\end{derivation}

\begin{derivation}[从\eqnrefs{eq:sigma-definition}到\eqnrefs{eq:sigma-sigma-trace-r68}]

正文\eqnrefs{eq:sigma-definition} 已给出 $\sigma^{\mu\nu}$ 与两个 $\gamma$ 矩阵对易子的关系。将两个这样的对象相乘并展开成四伽马迹后，所有张量结构都能由已知迹公式缩并，最终得到正文\eqnrefs{eq:sigma-sigma-trace-r68}。



定义
\begin{equation*}
\sigma^{\mu\nu}=\frac{i}{2}(\gamma^\mu\gamma^\nu-\gamma^\nu\gamma^\mu).
\end{equation*}
所以
\begin{align*}
\Tr(\sigma^{\mu\nu}\sigma^{\rho\sigma})
=-\frac14\Tr\Big[&
\gamma^\mu\gamma^\nu\gamma^\rho\gamma^\sigma
-\gamma^\mu\gamma^\nu\gamma^\sigma\gamma^\rho\\
&-\gamma^\nu\gamma^\mu\gamma^\rho\gamma^\sigma
+\gamma^\nu\gamma^\mu\gamma^\sigma\gamma^\rho
\Big].
\end{align*}
对四项分别代入
\(
\Tr(\gamma^a\gamma^b\gamma^c\gamma^d)
=4(\eta^{ab}\eta^{cd}-\eta^{ac}\eta^{bd}+\eta^{ad}\eta^{bc})
\)。四个迹依次为
\begin{align*}
T_1&=4(\eta^{\mu\nu}\eta^{\rho\sigma}
-\eta^{\mu\rho}\eta^{\nu\sigma}
+\eta^{\mu\sigma}\eta^{\nu\rho}),\\
T_2&=4(\eta^{\mu\nu}\eta^{\rho\sigma}
-\eta^{\mu\sigma}\eta^{\nu\rho}
+\eta^{\mu\rho}\eta^{\nu\sigma}),\\
T_3&=4(\eta^{\mu\nu}\eta^{\rho\sigma}
-\eta^{\nu\rho}\eta^{\mu\sigma}
+\eta^{\nu\sigma}\eta^{\mu\rho}),\\
T_4&=4(\eta^{\mu\nu}\eta^{\rho\sigma}
-\eta^{\nu\sigma}\eta^{\mu\rho}
+\eta^{\nu\rho}\eta^{\mu\sigma}).
\end{align*}
原组合是 $-(T_1-T_2-T_3+T_4)/4$。代入后，四个
$\eta^{\mu\nu}\eta^{\rho\sigma}$ 项的系数为 $1-1-1+1=0$；其余两种张量结构分别收集为
\begin{align*}
-\frac14(T_1-T_2-T_3+T_4)
&=4\eta^{\mu\rho}\eta^{\nu\sigma}
-4\eta^{\mu\sigma}\eta^{\nu\rho}.
\end{align*}
因此
\begin{equation*}
{
\Tr(\sigma^{\mu\nu}\sigma^{\rho\sigma})
=4(\eta^{\mu\rho}\eta^{\nu\sigma}
-\eta^{\mu\sigma}\eta^{\nu\rho}).}
\end{equation*}
进一步令 $\rho=\mu$、$\sigma=\nu$ 并求和，得到
\begin{align*}
\Tr(\sigma^{\mu\nu}\sigma_{\mu\nu})
&=4(\delta^\mu_\mu\delta^\nu_\nu-\delta^\mu_\nu\delta^\nu_\mu)\\
&=4(16-4)=48.
\end{align*}
由于 $\sigma^{\mu\nu}\sigma_{\mu\nu}$ 是 Lorentz 标量并在不可约 Dirac 矩阵空间中与恒等式成正比，而 $\Tr I_4=4$，因此
\begin{equation*}
\sigma^{\mu\nu}\sigma_{\mu\nu}=12I_4.
\end{equation*}
这两个式子是磁矩、Pauli 形状因子和张量算符中最常用的 Clifford 收缩。
\end{derivation}

\begin{derivation}[从\eqnrefs{eq:sigma-definition}到\eqnrefs{eq:sigma-gamma5-dual-r68}]
在四维和当前约定 $\epsilon^{0123}=+1$ 下，考虑两个秩二反对称张量

\begin{equation*}
\sigma^{\mu\nu}\gamma^5,
\qquad
\epsilon^{\mu\nu\rho\sigma}\sigma_{\rho\sigma}.
\end{equation*}
两者的自由指标都是 $\mu,\nu$，并且都在 Lorentz 变换下按反对称二阶张量变换，因此它们之间只能相差一个与指标无关的常数。把关系写成
\begin{equation*}
\sigma^{\mu\nu}\gamma^5
=C_D\epsilon^{\mu\nu\rho\sigma}\sigma_{\rho\sigma}.
\end{equation*}
取 $(\mu,\nu)=(0,1)$ 来求 $C_D$。由于
$\gamma^1\gamma^0=-\gamma^0\gamma^1$，有
\begin{align*}
\sigma^{01}
&=\frac{i}{2}(\gamma^0\gamma^1-\gamma^1\gamma^0),\\
&=i\gamma^0\gamma^1.
\end{align*}
所以
\begin{align*}
\sigma^{01}\gamma^5
&=(i\gamma^0\gamma^1)
(i\gamma^0\gamma^1\gamma^2\gamma^3),\\
&=-\gamma^0\gamma^1\gamma^0\gamma^1\gamma^2\gamma^3,\\
&=+(\gamma^0)^2(\gamma^1)^2\gamma^2\gamma^3,\\
&=-\gamma^2\gamma^3.
\end{align*}
另一方面，$\epsilon^{01\rho\sigma}$ 只有 $(\rho,\sigma)=(2,3)$ 和 $(3,2)$ 两项：
\begin{align*}
\epsilon^{01\rho\sigma}\sigma_{\rho\sigma}
&=\epsilon^{0123}\sigma_{23}+\epsilon^{0132}\sigma_{32},\\
&=\sigma_{23}+(-1)(-\sigma_{23}),\\
&=2\sigma_{23}.
\end{align*}
对空间指标降指标不改变 $\sigma_{23}$ 的符号，而
\begin{align*}
\sigma_{23}=\sigma^{23}
&=\frac{i}{2}(\gamma^2\gamma^3-\gamma^3\gamma^2),\\
&=i\gamma^2\gamma^3.
\end{align*}
因此右边等于 $2i\gamma^2\gamma^3$。比较
$-\gamma^2\gamma^3=C_D(2i\gamma^2\gamma^3)$ 得到
$C_D=i/2$，于是
\begin{equation*}
{
\sigma^{\mu\nu}\gamma^5
=\frac{i}{2}\epsilon^{\mu\nu\rho\sigma}\sigma_{\rho\sigma}.}
\end{equation*}
若翻转度规或 $\epsilon^{0123}$ 约定，右端符号必须同步改变。因此该对偶恒等式的符号由度规与 $\epsilon^{0123}$ 的取向约定共同固定。
\end{derivation}
